\documentclass{article}

\usepackage{amsmath}
\usepackage{amsthm}
\usepackage{amssymb}
\usepackage[hidelinks]{hyperref}

\newtheorem{lemma}{Lemma}
\newtheorem{theorem}{Theorem}
\providecommand*{\lemmaautorefname}{Lemma}

\newtheorem{corollary}[lemma]{Corollary}
\providecommand*{\corollaryautorefname}{Corollary}

\renewcommand{\refname}{\raggedright References}

\begin{document}

\begin{center}
\section*{Derivation of a Quantile Confidence Interval Based on Order Statistics}
{\small Frederik Danielsen}
\end{center}

\vspace{1em}

\noindent This note provides a mathematical derivation of a confidence interval for a quantile $\xi_q$ based on order statistics. The results presented here are classical and are not original to this note. They are derived independently to ensure a thorough understanding of the underlying mathematics, to present the arguments in a form suited to the Problab\footnote{\url{https://github.com/FrederikDanielsenDevelopment/problab}} project, and to obtain conclusions that are directly useful for the implementation in Problab. The derivation assumes i.i.d. samples $X_1,\ldots,X_n$ satisfying $P(X_i\leq\xi_q)=q$ and a significance level $\alpha\in(0,1)$, corresponding to a target coverage level of $1-\alpha$. In the first lemma, a confidence interval of the form $[X_{(k_L^*)},X_{(k_U^*+1)}]$ is derived by choosing the indices $k_L^*$ and $k_U^*$ such that the two binomial tail probabilities are bounded by $\alpha/2$. In the second lemma, it is shown that the iterative procedures used in the implementation terminate at exactly these optimal indices.

Practically, these results enable Problab to return finite-sample confidence intervals with coverage probability at least $1-\alpha$ for estimated quantiles, such as medians or percentiles.

\vspace{2em}


\begin{lemma}\label{lemma:quantile_CI}
Let $X_i$ with $i\in\{1,\ldots,n-1\}$ be i.i.d samples. Suppose that $P(X_i\leq \xi_q)=q$ for all $i$, some $q\in(0,1)$, and $\xi_q$ in the value space of $X_i$. Let
$$I_i =
\begin{cases}
1,\quad X_i\leq\xi_q,\\
0,\quad X_i>\xi_q.
\end{cases}$$
and let
$$K=\sum_{i=1}^n I_i\sim Bin(n,q).$$
Let $k_L^*$ and $k_U^*$ solve
$$
\begin{aligned}
k_L^*
&=
\max_{k_L\in\{1,\ldots,n-1\}}
\left\{
k_L:
P(K<k_L)\leq\frac{\alpha}{2}
\right\}\\
k_U^*
&=
\min_{k_U\in\{1,\ldots,n-1\}}
\left\{
k_U:
P(K>k_U)\leq\frac{\alpha}{2}
\right\}.
\end{aligned}
$$
Then
$$[X_{(k_L^*)},X_{(k_U^*+1)}]$$
is a confidence interval for $\xi_q$ with coverage probability at least $1-\alpha$. Furthermore, among the intervals satisfying the two tail constraints, solving the optimization problems above produces the narrowest interval.
\end{lemma}
\begin{proof}
$I_i$ can be regarded as the outcomes of individual Bernouli trials with probability of success $q$ with
$$K=\sum_{i=1}^nI_i\sim Bin(n,q)$$
denoting the number of samples with values at most $\xi_q$. Now, consider the ordered samples
$$X_{(1)},X_{(2)},\ldots,X_{(n)}.$$
If $K\leq k_U$ for some $k_U\in\{1,\ldots,n-1\}$, that would imply that at most $k_U$ samples fell below $\xi_q$. That is, $\xi_q<X_{(k_U+1)}$. Likewise, if $K\geq k_L$ for some $k_L\in\{1,\ldots,n-1\}$, that would imply that at least $k_L$ samples fell below $\xi_q$. That is, $\xi_q\geq X_{(k_L)}$. Hence,
$$P(k_L\leq K\leq k_U)=P(X_{(k_L)}\leq\xi_q<X_{(k_U+1)}).$$
Therefore, if we determine $k_L$ and $k_U$ such that
$$P(k_L\leq K \leq k_U)=1-\alpha,$$
for some $\alpha\in(0,1)$, then the interval $[X_{(k_L)},X_{(k_U+1)}]$ is a confidence interval for $\xi_q$ with coverage probability $1-\alpha$. However, this equality is one with two unknowns and thus multiple solutions. However, we may impose an additional constraint, namely that the likelihood that $\xi_q$ lies above the confidence interval is the same as the likelihood that it lies below it $-$ at least to the extent that a discrete number of samples allows it. Specifically, we require that
$$P(\xi_q\geq X_{(k_U+1)})=P(\xi_q<X_{(k_L)})=\alpha/2.$$
However, as mentioned, since we are dealing with a discrete number of samples, it is likely that no $(k_L,k_U)$ exists that fulfil this exactly. Therefore, we make do with the best estimate (likely more conservative than $1-\alpha$) and require simply that
$$
\begin{aligned}
P(\xi_q \geq X_{(k_U+1)})&\leq \alpha/2,\\
P(\xi_q<X_{(k_L)})&\leq \alpha/2\\
\Leftrightarrow\\
P(K>k_U)&\leq \alpha/2,\\
P(K<k_L)&\leq \alpha/2.
\end{aligned}
$$
Since $\alpha<1$, any pair $(k_L,k_U)$ satisfying both tail constraints necessarily satisfies $k_L\leq k_U$. Indeed, if $k_L>k_U$, then every possible value of $K$ satisfies either $K<k_L$ or $K>k_U$, and therefore
$$
1=P\big((K<k_L)\cup(K>k_U)\big)
\leq P(K<k_L)+P(K>k_U)
\leq \alpha,
$$
which is a contradiction. Furthermore, in case you were wondering, since the lower endpoint $X_{(k_L)}$ requires $k_L\geq1$ and the upper endpoint $X_{(k_U+1)}$ requires $k_U\leq n-1$, the inequality $k_L\leq k_U$ implies that any feasible pair must satisfy
$$
k_L,k_U\in\{1,\ldots,n-1\}.
$$
To obtain the narrowest confidence interval satisfying these constraints, we choose the largest possible $k_L$ and the smallest possible $k_U$. Thus, we solve
$$
\begin{aligned}
k_L^*
&=
\max_{k_L\in\{1,\ldots,n-1\}}
\left\{
k_L:
P(K<k_L)\leq\frac{\alpha}{2}
\right\},\\
k_U^*
&=
\min_{k_U\in\{1,\ldots,n-1\}}
\left\{
k_U:
P(K>k_U)\leq\frac{\alpha}{2}
\right\}.
\end{aligned}
$$
In this way we ensure that
$$P(k_L^*\leq K \leq k_U^*)=P(X_{(k_L^*)}\leq \xi_q< X_{(k_U^*+1)})\geq 1-\alpha.$$
Then the interval
$$[X_{(k_L^*)},X_{(k_U^*+1)}]$$
is at least as conservative as a confidence interval with coverage probability $1-\alpha$.
\end{proof}

\begin{lemma}
Let $\alpha\in(0,1)$, $q\in(0,1)$, $n\geq\lceil\ln(\alpha/2)/\ln(\max (q, 1-q))\rceil$ and $K\sim Bin(n,q)$. Then the optimization problems
$$
\begin{aligned}
k_L^*
&=
\max_{k\in\{1,\ldots,n-1\}}
\left\{
k:
P(K<k)\leq\frac{\alpha}{2}
\right\},\\
k_U^*
&=
\min_{k\in\{1,\ldots,n-1\}}
\left\{
k:
P(K>k)\leq\frac{\alpha}{2}
\right\}.
\end{aligned}
$$
have solutions.
\end{lemma}
\begin{proof}
Suppose that no feasible pair $(k_L^*, k_U^*)$ exists. Since $P(K<k)$ is increasing and $P(K>k)$ is decreasing for $k\in\{1,\ldots,n-1\}$, that would imply that
$$
P(K<1)=P(K=0)=(1-q)^n>\alpha/2
\,\bigvee\,
P(K>n-1)=P(K=n)=q^n>\alpha/2
$$
$$
\Leftrightarrow
\max(q^n,(1-q)^n)=\max(q,1-q)^n>\alpha/2.
$$
Therefore, if we want
$$
P(K<1)\leq\alpha/2
\,\wedge\,
P(K>n-1)\leq\alpha/2,
$$
then it must be the case that
$$
\max(q,1-q)^n\leq \frac{\alpha}{2}.
$$
Taking logarithms gives
$$
n\ln(\max(q,1-q))
\leq
\ln\left(\frac{\alpha}{2}\right).
$$
Since $\max(q,1-q)\in(0,1)$, we have
$$
\ln(\max(q,1-q))<0,
$$
so dividing by $\ln(\max(q,1-q))$ reverses the inequality:
$$
n \geq \frac{\ln(\alpha/2)}{\ln(\max(q,1-q))}.
$$
Thus, the minimum value required for $n$ is
$$
n_\mathrm{min}
=
\left\lceil
\frac{\ln(\alpha/2)}
{\ln(\max(q,1-q))}
\right\rceil.
$$
Hence, under the stated condition on $n$, both optimization problems have solutions.
\end{proof}




\begin{lemma}
Let $K\sim Bin(n,q)$ and $\alpha\in(0,1)$, and suppose that the optimization problems
$$
\begin{aligned}
k_L^*
&=
\max_{k\in\{1,\ldots,n-1\}}
\left\{
k:
P(K<k)\leq\frac{\alpha}{2}
\right\},\\
k_U^*
&=
\min_{k\in\{1,\ldots,n-1\}}
\left\{
k:
P(K>k)\leq\frac{\alpha}{2}
\right\}
\end{aligned}
$$
have solutions. Consider the binary-search procedures
$$
\begin{aligned}
&a_L\leftarrow 1,\qquad b_L\leftarrow n-1,\\
&\text{while } a_L<b_L:\\
&\qquad m_L\leftarrow
\left\lceil\frac{a_L+b_L}{2}\right\rceil,\\
&\qquad\text{if }
P(K<m_L)\leq\frac{\alpha}{2},
\text{ then }
a_L\leftarrow m_L,\\
&\qquad\text{otherwise }
b_L\leftarrow m_L-1,
\end{aligned}
$$
and
$$
\begin{aligned}
&a_U\leftarrow 1,\qquad b_U\leftarrow n-1,\\
&\text{while } a_U<b_U:\\
&\qquad m_U\leftarrow
\left\lfloor\frac{a_U+b_U}{2}\right\rfloor,\\
&\qquad\text{if }
P(K>m_U)\leq\frac{\alpha}{2},
\text{ then }
b_U\leftarrow m_U,\\
&\qquad\text{otherwise }
a_U\leftarrow m_U+1.
\end{aligned}
$$
Then both procedures terminate with
$$
a_L=b_L=k_L^*
\qquad\text{and}\qquad
a_U=b_U=k_U^*.
$$
Each procedure evaluates its tail probability at most
$$
\left\lceil\log_2(n-1)\right\rceil
$$
times.
\end{lemma}

\begin{proof}
For the lower index, $P(K<k)$ is nondecreasing in $k$. Since $k_L^*$ exists, $1$ is feasible, and every feasible index is at most $k_L^*$. Thus, initially,
$$
a_L\leq k_L^*\leq b_L.
$$
At each iteration, $m_L$ is a candidate index. If
$$
P(K<m_L)\leq\frac{\alpha}{2},
$$
then $m_L$ is feasible, so $k_L^*\geq m_L$ and replacing $a_L$ by $m_L$ preserves
$$
a_L\leq k_L^*\leq b_L.
$$
Otherwise, $m_L$ is infeasible. Since the lower-tail probability is nondecreasing, every index at least $m_L$ is infeasible, so $k_L^*\leq m_L-1$. Replacing $b_L$ by $m_L-1$ therefore preserves the same inequality.
Because $m_L$ is rounded up, each iteration strictly reduces the number of candidate indices. The procedure terminates when $a_L=b_L$. At that point, the preserved inequality gives
$$
a_L=b_L=k_L^*.
$$
For the upper index, $P(K>k)$ is nonincreasing in $k$. Since $k_U^*$ exists, $n-1$ is feasible, and every feasible index is at least $k_U^*$. Thus, initially,
$$
a_U\leq k_U^*\leq b_U.
$$
If
$$
P(K>m_U)\leq\frac{\alpha}{2},
$$
then $m_U$ is feasible, so $k_U^*\leq m_U$. Replacing $b_U$ by $m_U$ preserves
$$
a_U\leq k_U^*\leq b_U.
$$
Otherwise, $m_U$ is infeasible. Since the upper-tail probability is nonincreasing, every index at most $m_U$ is infeasible, so $k_U^*\geq m_U+1$. Replacing $a_U$ by $m_U+1$ therefore preserves the same inequality.
Because $m_U$ is rounded down, each iteration strictly reduces the number of candidate indices. The procedure terminates when $a_U=b_U$, and the preserved inequality gives
$$
a_U=b_U=k_U^*.
$$
Each iteration removes at least half of the remaining candidate indices. Since there are initially $n-1$ candidates, each procedure requires at most
$$
\left\lceil\log_2(n-1)\right\rceil
$$
tail-probability evaluations.
\end{proof}



\end{document}

